% SYNC-ID: problem-definition
\section{Problem Definition}\label{sec:problem-definition}
% CONTENT_DEFERRED_TO_WRITING_PHASE
\textit{(THIS IS PLACEHOLDER TEXT, PLZ FULFILL IT WHEN WRITING.)}
% SYNC-ID: problem-placeholder-1
The problem statement will be written after the scope and terminology have been approved. This fixture shows the expected paragraph rhythm while making no operational definition or research question.
% SYNC-ID: problem-placeholder-2
Notation can be introduced in this position once the paired definitions are frozen. No symbol in the current source carries a paper-specific meaning.
% SYNC-ID: problem-placeholder-3
This paragraph reserves an input and output description without filling in values, examples, or evaluation conditions.
% SYNC-ID: problem-placeholder-4
The final terminology pass will connect each term to the approved glossary and its Chinese counterpart.
% SYNC-ID: problem-placeholder-5
For layout purposes, this section reserves a final bridge from terminology to the later study sections. The bridge should describe only approved interfaces and should preserve the same qualifiers in both languages.
% SYNC-ID: problem-placeholder-6
The placeholder further reserves room for a compact formal statement and a plain-language paraphrase. Both versions will be generated from one reviewed semantic block, so neither can silently narrow or broaden the other.
